\documentclass[12pt]{article}
\usepackage[ukrainian]{babel}
\usepackage{fontspec}
% --- НАЛАШТУВАННЯ ШРИФТІВ ---
\setmainfont{Schoolbook-Regular.otf}[
    Path = ../,
    BoldFont = Schoolbook-Bold.otf,
    ItalicFont = Schoolbook-Italic.otf,
    BoldItalicFont = Schoolbook-BoldItalic.otf
]
\usepackage[italic]{mathastext}
\usepackage[a4paper,margin=1.8cm,bottom=2cm,top=2cm]{geometry}
\usepackage{amsmath,amssymb}
\usepackage{tikz}
\usepackage{pgfplots}
\pgfplotsset{compat=1.16}
\usetikzlibrary{calc,patterns,angles,quotes,intersections,babel}
\usepackage{xcolor,array,fancyhdr,enumitem,multicol}
\usepackage{colortbl}
\usepackage{tcolorbox}
\tcbuselibrary{skins,breakable}
\usepackage[hidelinks]{hyperref}
\usepackage[firstpage=false, text=@pvtr2525, scale=0.6, color=gray!12]{draftwatermark}

% Заборона розриву інлайн-формул
\binoppenalty=10000
\relpenalty=10000

\definecolor{mainGreen}{RGB}{34, 120, 64}
\definecolor{yearOrange}{RGB}{220, 100, 30}
\definecolor{titleBg}{RGB}{225, 240, 220}
\definecolor{instrBg}{RGB}{255, 240, 220}
\definecolor{instrBorder}{RGB}{220, 150, 80}

\pagestyle{fancy}
\fancyhf{}
\renewcommand{\headrulewidth}{0.4pt}
\renewcommand{\footrulewidth}{0.4pt}
\fancyhead[C]{\small\color{gray!80} Збірник НМТ \textendash{} Варіант за 9 червня}
\fancyhead[R]{\small\color{mainGreen}\textbf{\href{https://t.me/pvtr2525}{@pvtr2525}}}
\fancyfoot[L]{\small\color{gray!80}\href{https://t.me/pvtr2525}{ТГ: @pvtr2525}}
\fancyfoot[C]{\small\color{gray!80}\thepage}
\fancyfoot[R]{\small\color{gray!80}НМТ 2026}

\newcommand{\answerTable}[5]{%
\par\vspace{0.15cm}
\noindent
\begin{tabular}{|*{5}{>{\centering\arraybackslash}m{3cm}|}}
\hline
\rule[-0.15cm]{0pt}{0.6cm}\textbf{А} & \rule[-0.15cm]{0pt}{0.6cm}\textbf{Б} & \rule[-0.15cm]{0pt}{0.6cm}\textbf{В} & \rule[-0.15cm]{0pt}{0.6cm}\textbf{Г} & \rule[-0.15cm]{0pt}{0.6cm}\textbf{Д} \\
\hline
\rule[-0.2cm]{0pt}{0.75cm}#1 & \rule[-0.2cm]{0pt}{0.75cm}#2 & \rule[-0.2cm]{0pt}{0.75cm}#3 & \rule[-0.2cm]{0pt}{0.75cm}#4 & \rule[-0.2cm]{0pt}{0.75cm}#5 \\
\hline
\end{tabular}
\par\vspace{0.25cm}
}

\newcommand{\answerTableTall}[5]{%
\par\vspace{0.15cm}
\noindent
\begin{tabular}{|*{5}{>{\centering\arraybackslash}m{3cm}|}}
\hline
\rule[-0.2cm]{0pt}{0.7cm}\textbf{А} & \rule[-0.2cm]{0pt}{0.7cm}\textbf{Б} & \rule[-0.2cm]{0pt}{0.7cm}\textbf{В} & \rule[-0.2cm]{0pt}{0.7cm}\textbf{Г} & \rule[-0.2cm]{0pt}{0.7cm}\textbf{Д} \\
\hline
\rule[-0.45cm]{0pt}{1.2cm}#1 & \rule[-0.45cm]{0pt}{1.2cm}#2 & \rule[-0.45cm]{0pt}{1.2cm}#3 & \rule[-0.45cm]{0pt}{1.2cm}#4 & \rule[-0.45cm]{0pt}{1.2cm}#5 \\
\hline
\end{tabular}
\par\vspace{0.25cm}
}

\newcommand{\instructionBox}[1]{%
\par\vspace{0.3cm}
\begin{tcolorbox}[colback=instrBg, colframe=instrBorder, boxrule=0.6pt, arc=2pt,
    left=8pt, right=8pt, top=4pt, bottom=4pt]
\centering\small\bfseries #1
\end{tcolorbox}
\par\vspace{0.15cm}
}

\newcommand{\sectionTitle}[1]{%
\par\vspace{0.4cm}
\begin{tcolorbox}[colback=titleBg, colframe=titleBg, boxrule=0pt, arc=8pt, halign=center, fontupper=\Large\bfseries\color{mainGreen}]
#1
\end{tcolorbox}\par\vspace{0.3cm}}
\newcommand{\chapterTitle}[1]{\clearpage\sectionTitle{#1}}

\newcommand{\nmtAnswerBox}{%
\par\vspace{0.2cm}\noindent
Відповідь:\ \ 
\begin{tabular}{@{}*{4}{|>{\centering\arraybackslash}p{0.8cm}}|@{\hspace{0.1cm}\textbf{,}\hspace{0.1cm}}*{3}{|>{\centering\arraybackslash}p{0.8cm}}|@{}}
\hline
\rule[-0.2cm]{0pt}{0.7cm} & & & & & & \\
\hline
\end{tabular}
\par\vspace{0.3cm}}

\newcommand{\matchingGrid}{%
    \begingroup
    \renewcommand{\arraystretch}{1.15}
    \setlength{\tabcolsep}{4pt}
    \footnotesize
    \begin{tabular}{r|c|c|c|c|c|}
         \multicolumn{1}{c}{} & \multicolumn{1}{c}{\textbf{А}} & \multicolumn{1}{c}{\textbf{Б}} & \multicolumn{1}{c}{\textbf{В}} & \multicolumn{1}{c}{\textbf{Г}} & \multicolumn{1}{c}{\textbf{Д}} \\ \cline{2-6}
         \textbf{1} & \phantom{X} & \phantom{X} & \phantom{X} & \phantom{X} & \phantom{X} \\ \cline{2-6}
         \textbf{2} & & & & & \\ \cline{2-6}
         \textbf{3} & & & & & \\ \cline{2-6}
    \end{tabular}
    \endgroup
}

\newcounter{zad}
\newcommand{\zadtask}[1]{\stepcounter{zad}\par\vspace{0.3cm}\noindent\textbf{\thezad.}\ \ #1\par\vspace{0.2cm}}
\newcommand{\zadnum}{\stepcounter{zad}\noindent\textbf{\thezad.}\ \ }

\begin{document}
\thispagestyle{empty}
\vspace*{1.2cm}
\begin{center}
{\Large\bfseries\color{mainGreen} РЕАЛЬНИЙ ВАРІАНТ НМТ-2026}\\[0.4cm]
{\Huge\bfseries\color{mainGreen} ВАРІАНТ ЗА 9 ЧЕРВНЯ}\\[0.6cm]
{\large Real Test Notebook з усіма геометричними рисунками}\\[1cm]
{\large\color{mainGreen}\textbf{\href{https://t.me/pvtr2525}{Телеграм: @pvtr2525}}}
\end{center}
\clearpage

\chapterTitle{ТЕСТОВИЙ ЗОШИТ (9 ЧЕРВНЯ)}
\instructionBox{Завдання 1--15 мають по п'ять варіантів відповіді, з яких лише один правильний. Виберіть правильний варіант відповіді.}

\zadtask{Знайдіть найменший цілий розв'язок нерівності $2x+7>27$.}
\answerTable{$20$}{$15$}{$11$}{$10$}{$9$}

\noindent
\begin{minipage}[c]{0.42\textwidth}
\zadnum На діаграмі відображено дані про швидкість автомобілів, яку зафіксували камери спостереження у вечірній час. За даними діаграми визначте загальну кількість автомобілів, рух яких було зафіксовано.
\end{minipage}%
\hfill
\begin{minipage}[c]{0.56\textwidth}
\centering
\begin{tikzpicture}[scale=0.9, thick, line join=round]
    \draw[->, gray!80] (0,0) -- (0,5.2) node[above, rotate=90, anchor=south, xshift=-2.2cm, yshift=0.5cm, text=gray!80!black, font=\scriptsize] {Кількість автомобілів};
    \draw[->, gray!80] (0,0) -- (7.5,0) node[below, text=gray!80!black, font=\scriptsize, xshift=-3.5cm, yshift=-0.5cm] {Швидкість автомобіля, км/год};
    
    \foreach \y/\val in {0.6/2, 1.2/4, 1.8/6, 2.4/8, 3.0/10, 3.6/12, 4.2/14, 4.8/16} {
        \draw[gray!30, thin] (0,\y) -- (7,\y);
        \node[left, font=\scriptsize, text=gray!80!black] at (0,\y) {\val};
    }
    \node[left, font=\scriptsize, text=gray!80!black] at (0,0) {0};
    
    \filldraw[fill=cyan!60!blue, draw=none] (0.5,0) rectangle (1.2, 2.4); 
    \filldraw[fill=cyan!60!blue, draw=none] (1.9,0) rectangle (2.6, 3.6); 
    \filldraw[fill=cyan!60!blue, draw=none] (3.3,0) rectangle (4.0, 4.5); 
    \filldraw[fill=cyan!60!blue, draw=none] (4.7,0) rectangle (5.4, 3.0); 
    \filldraw[fill=cyan!60!blue, draw=none] (6.1,0) rectangle (6.8, 1.2); 
    
    \node[below, font=\scriptsize] at (0.85,0) {40--60};
    \node[below, font=\scriptsize] at (2.25,0) {61--80};
    \node[below, font=\scriptsize] at (3.65,0) {81--100};
    \node[below, font=\scriptsize] at (5.05,0) {101--120};
    \node[below, font=\scriptsize] at (6.45,0) {понад 120};
\end{tikzpicture}
\end{minipage}
\par\vspace{0.2cm}
\answerTable{$48$}{$49$}{$50$}{$45$}{$52$}

\zadtask{Дарина купила сир та фрукти, витративши $240$~грн. Скільки грошей Дарина витратила на фрукти, якщо за сир вона заплатила $\dfrac{3}{5}$ витраченої суми?}
\answerTable{$34$~грн}{$30$~грн}{$45$~грн}{$60$~грн}{$96$~грн}

\noindent
\begin{minipage}[c]{0.56\textwidth}
\zadnum На рисунку 1 схематично зображено бічну проєкцію будівельної драбини-трансформера (рис. 2), розкладеної у вигляді помосту на горизонтальній підлозі (пряма $AD$). Чотирикутник $ABCD$~--- рівнобічна трапеція. Визначте градусну міру кута $BAD$, якщо $\angle BCD = 110^\circ$.
\end{minipage}%
\hfill
\begin{minipage}[c]{0.42\textwidth}
\centering
\begin{tikzpicture}[scale=0.7, thick, line join=round]
    \coordinate (A) at (0,0);
    \coordinate (D) at (5,0);
    \coordinate (B) at (1.5,2.5);
    \coordinate (C) at (3.5,2.5);
    \draw (A) -- (B) -- (C) -- (D) -- cycle;
    \draw[gray!80, thin] (-0.5,0) -- (5.5,0);
    \node[below left] at (A) {$A$};
    \node[above left] at (B) {$B$};
    \node[above right] at (C) {$C$};
    \node[below right] at (D) {$D$};
    \pic[draw, angle radius=0.4cm, pic text={$110^\circ$}, angle eccentricity=1.7, font=\scriptsize] {angle = B--C--D};
    \pic[draw, angle radius=0.4cm, pic text={$?$}, angle eccentricity=1.5, font=\scriptsize] {angle = D--A--B};
    \node[below] at (2.5, -0.5) {Рис. 1};
\end{tikzpicture}
\end{minipage}
\par\vspace{0.2cm}
\answerTable{$70^\circ$}{$55^\circ$}{$60^\circ$}{$110^\circ$}{$35^\circ$}

\zadtask{Куля утворена обертанням
\vspace{0.1cm}
\begin{enumerate}[label=\textbf{\Alph*}, leftmargin=2.8em, itemsep=0.1cm]
    \item кола навколо його діаметра
    \item круга навколо прямої, що проходить через його центр
    \item кола навколо його хорди
    \item круга навколо прямої, що не проходить через його центр
    \item кола навколо дотичної
\end{enumerate}
\vspace{0.1cm}}

\zadtask{$|5^{-1} - 3| =$}
\answerTable{$-2{,}8$}{$2{,}8$}{$-8$}{$2$}{$8$}

\zadtask{Розв'яжіть рівняння $3^{x} = 81 \cdot 9^{x}$.}
\answerTable{$-2$}{$-4$}{$-3$}{$3$}{$4$}

\noindent
\begin{minipage}[c]{0.56\textwidth}
\zadnum На рисунку зображено графік функції $y=f(x)$, визначеної на проміжку $[-4;\, 4]$. Укажіть проміжок, на якому функція набуває лише від'ємних значень.
\end{minipage}%
\hfill
\begin{minipage}[c]{0.42\textwidth}
\centering
\begin{tikzpicture}[scale=0.6, thick]
    \draw[step=1cm, gray!40, thin] (-4,-2) grid (5,3);
    \draw[->] (-4.5,0) -- (5.5,0) node[right] {$x$};
    \draw[->] (0,-2.5) -- (0,3.5) node[above] {$y$};
    \node[below left] at (0,0) {$0$};
    \node[below right] at (1,0) {$1$};
    \node[above left] at (0,1) {$1$};
    \node[below] at (-4,0) {$-4$};
    \node[below] at (4,0) {$4$};
    \draw[magenta!80!black, thick, smooth] plot coordinates {(-4,-1) (-2,1.5) (0,1) (1,0) (2,-1) (3,0) (4,1)};
    \fill (-4,-1) circle (3pt);
    \fill (4,1) circle (3pt);
    \node[above right] at (1,1.5) {$y=f(x)$};
\end{tikzpicture}
\end{minipage}
\par\vspace{0.2cm}
\answerTable{$(-4; -1)$}{$(0; 1)$}{$(1; 2)$}{$(1; 4)$}{$(3; 4)$}

\noindent
\begin{minipage}[c]{0.62\textwidth}
\zadnum Трикутник $ABC$ вписано в коло з центром у точці $O$, що належить стороні $AC$ (див. рисунок). Які з наведених тверджень є правильними?
\vspace{0.1cm}
\begin{enumerate}[label=\textbf{\Roman*.}, leftmargin=2.8em, itemsep=0.1cm]
\item $\angle ABC$ є прямим.
\item $\angle BAC$ вдвічі менший від $\angle BOC$.
\item Довжина $AC$ вдвічі більша за довжину $BO$.
\end{enumerate}
\end{minipage}%
\hfill
\begin{minipage}[c]{0.36\textwidth}
\centering
\begin{tikzpicture}[scale=1, thick, line join=round]
    \coordinate (O) at (0,0);
    \draw (O) circle (2);
    \coordinate (A) at (-1.732, -1);
    \coordinate (C) at (1.732, 1);
    \coordinate (B) at (-0.5, 1.936);
    \draw (A) -- (B) -- (C) -- cycle;
    \draw (A) -- (C);
    \draw (B) -- (O);
    \fill (O) circle (2pt) node[below right] {$O$};
    \node[below left] at (A) {$A$};
    \node[above left] at (B) {$B$};
    \node[above right] at (C) {$C$};
\end{tikzpicture}
\end{minipage}
\par\vspace{0.2cm}
\answerTable{лише I}{лише I та II}{лише I та III}{лише II та III}{I, II та III}

\noindent
\begin{minipage}[c]{0.62\textwidth}
\zadnum Обчисліть площу $S$ заштрихованої фігури, обмеженої графіком функції $y=e^x$, осями координат і прямою $x=1$.
\vspace{0.2cm}
\par\noindent\textbf{А}\ \ $S = e+1$
\par\noindent\textbf{Б}\ \ $S = e$
\par\noindent\textbf{В}\ \ $S = 1$
\par\noindent\textbf{Г}\ \ $S = 1-e$
\par\noindent\textbf{Д}\ \ $S = e-1$
\end{minipage}%
\hfill
\begin{minipage}[c]{0.36\textwidth}
\centering
\begin{tikzpicture}[scale=1.5, thick]
    \draw[->] (-1.5,0) -- (2,0) node[right] {$x$};
    \draw[->] (0,-0.5) -- (0,3) node[above] {$y$};
    \draw[domain=-1.5:1.5, smooth, variable=\x, black] plot ({\x}, {exp(\x)});
    \fill[cyan!40] (0,0) -- (1,0) -- (1,{exp(1)}) -- plot[domain=1:0, smooth] ({\x}, {exp(\x)}) -- cycle;
    \draw (1,0) -- (1,{exp(1)});
    \node[below left] at (0,0) {$0$};
    \node[below] at (1,0) {$1$};
    \node[left] at (0,1) {$1$};
    \node[right] at (1.2, 2.5) {$y=e^x$};
\end{tikzpicture}
\end{minipage}
\par\vspace{0.3cm}

\zadtask{На столі лежать $18$ світлин, серед яких $8$ із зображенням природи, $6$ із зображенням людей та $4$ із зображенням будівель. Яка ймовірність того, що навмання вибрана світлина буде із зображенням природи або людини?}
\answerTableTall{$\dfrac{2}{9}$}{$\dfrac{4}{9}$}{$\dfrac{5}{9}$}{$\dfrac{7}{9}$}{$\dfrac{11}{18}$}

\noindent
\begin{minipage}[c]{0.56\textwidth}
\zadnum У прямокутній системі координат у просторі зображено прямокутний паралелепіпед $ABCDA_1B_1C_1D_1$, вершина $A$ якого збігається з початком координат, а ребра $AB$, $AD$ й $AA_1$ лежать на координатних осях (див. рисунок). Визначте координати вектора $\overrightarrow{DB_1}$.
\end{minipage}%
\hfill
\begin{minipage}[c]{0.42\textwidth}
\centering
\begin{tikzpicture}[scale=0.75, thick, line join=round]
    \coordinate (A) at (0,0);
    \coordinate (B) at (225:1.4);
    \coordinate (D) at (3.5,0);
    \coordinate (C) at ($(B)+(D)$);
    \coordinate (A1) at (0,4);
    \coordinate (B1) at ($(B)+(A1)$);
    \coordinate (C1) at ($(C)+(A1)$);
    \coordinate (D1) at ($(D)+(A1)$);
    
    % Axes
    \draw[->, >=stealth] (A) -- ($(B)+(225:1)$) node[left] {$x$};
    \draw[->, >=stealth] (A) -- ($(D)+(1,0)$) node[below] {$y$};
    \draw[->, >=stealth] (A) -- ($(A1)+(0,1.2)$) node[left] {$z$};
    
    % Edges
    \draw (B) -- (C) -- (D) -- (D1) -- (C1) -- (B1) -- cycle;
    \draw (C) -- (C1);
    \draw (A1) -- (B1);
    \draw (A1) -- (D1);
    \draw[dashed] (A) -- (B);
    \draw[dashed] (A) -- (D);
    \draw[dashed] (A) -- (A1);
    
    % Vector red dashed
    \draw[dashed, red, ->, >=stealth, line width=1.1pt] (D) -- (B1);
    
    % Text and ticks
    \node[left, xshift=-0.1cm] at (A) {$A$};
    \node[left] at (B) {$B$};
    \node[below, yshift=-0.1cm] at (B) {$2$};
    \node[below right] at (C) {$C$};
    \node[below, yshift=-0.1cm] at (D) {$3$};
    \node[above right] at (D1) {$D_1$};
    \node[above right] at (C1) {$C_1$};
    \node[above left] at (B1) {$B_1$};
    \node[left] at (A1) {$A_1$};
    \node[right, xshift=0.05cm] at (A1) {$4$};
\end{tikzpicture}
\end{minipage}
\par\vspace{0.2cm}
\answerTable{$(-2; 3; -4)$}{$(2; 3; 4)$}{$(2; -3; 4)$}{$(-2; -3; -4)$}{$(-2; 3; 4)$}

\zadtask{Спростіть вираз $25b - (b - 5)(5 + b) + 25$.}
\answerTableTall{$-b^2 + 25b + 50$}{$b^2 + 25b + 50$}{$-b^2 + 25$}{$-b^2 + 25b$}{$b^2 + 25b$}

\noindent
\begin{minipage}[c]{0.60\textwidth}
\zadnum На рисунку зображено рівнобедрений трикутник $ABC$, у якому $\angle ABC = 120^\circ$. Точки $K$ і $M$ є серединами сторін $AB$ та $AC$ відповідно. Визначте площу чотирикутника $KBCM$, якщо $AC = 24$ см.
\end{minipage}%
\hfill
\begin{minipage}[c]{0.38\textwidth}
\centering
\begin{tikzpicture}[scale=0.6, thick, line join=round]
    \coordinate (B) at (0,0);
    \coordinate (C) at (6,0);
    \coordinate (A) at (-3, 5.196);
    \coordinate (K) at (-1.5, 2.598);
    \coordinate (M) at (1.5, 2.598);
    
    \draw (A) -- (B) -- (C) -- cycle;
    \draw (K) -- (M);
    
    \node[above left] at (A) {$A$};
    \node[below left] at (B) {$B$};
    \node[below right] at (C) {$C$};
    \node[left] at (K) {$K$};
    \node[above right] at (M) {$M$};
    
    \pic[draw, angle radius=0.5cm, pic text={$120^\circ$}, angle eccentricity=1.5, font=\scriptsize] {angle = C--B--A};
\end{tikzpicture}
\end{minipage}
\par\vspace{0.2cm}
\answerTable{$24$~см$^2$}{$24\sqrt{3}$~см$^2$}{$72$~см$^2$}{$72\sqrt{3}$~см$^2$}{$36\sqrt{3}$~см$^2$}

\zadtask{Розв'яжіть систему рівнянь 
\[
\begin{cases} x + y = -4, \\ xy + y^2 = 12. \end{cases}
\] 
Якщо $(x_0; y_0)$~--- розв'язок системи, то $x_0 \cdot y_0 = $}
\answerTable{$-1$}{$3$}{$2$}{$-4$}{$-3$}
\newpage
\instructionBox{У завданнях 16--18 до кожного з трьох рядків інформації, позначених цифрами, доберіть один правильний варіант, позначений буквою.}

\zadtask{Узгодьте вираз (1--3) із проміжком (А--Д), якому належить значення цього виразу.
\vspace{0.3cm}
\noindent
\begin{minipage}[t]{0.45\textwidth}
\textit{Вираз}\par\vspace{0.15cm}
\textbf{1} \quad $\dfrac{\sqrt{200}}{\sqrt{8}}$\\[0.25cm]
\textbf{2} \quad $\log_8 200$\\[0.25cm]
\textbf{3} \quad $8 \cos 120^\circ$
\end{minipage}
\hfill
\begin{minipage}[t]{0.3\textwidth}
\textit{Проміжок}\par\vspace{0.15cm}
\textbf{А} \quad $[-5; 0)$\\[0.15cm]
\textbf{Б} \quad $[0; 2)$\\[0.15cm]
\textbf{В} \quad $[2; 4)$\\[0.15cm]
\textbf{Г} \quad $[4; 6)$\\[0.15cm]
\textbf{Д} \quad $[6; +\infty)$
\end{minipage}
\hfill
\begin{minipage}[t]{0.2\textwidth}
\vspace{0.4cm}
\begin{flushright}\matchingGrid\end{flushright}
\end{minipage}}

\zadtask{Доберіть до функції (1--3) її властивість (А--Д).\par
\vspace{0.3cm}
\noindent
\begin{minipage}[t]{0.30\textwidth}
\textit{Функція}\par\vspace{0.2cm}
\textbf{1} \quad $y = \sin x$\\[0.25cm]
\textbf{2} \quad $y = x^2 + 2$\\[0.25cm]
\textbf{3} \quad $y = 3x + 2$
\par\vspace{0.4cm}
\matchingGrid
\end{minipage}%
\hfill
\begin{minipage}[t]{0.66\textwidth}
\textit{Властивість функції}\par\vspace{0.2cm}
\textbf{А} \quad парна\\[0.15cm]
\textbf{Б} \quad графік функції симетричний відносно осі $x$\\[0.15cm]
\textbf{В} \quad спадає на області визначення\\[0.15cm]
\textbf{Г} \quad зростає на області визначення\\[0.15cm]
\textbf{Д} \quad графік функції проходить через точку $(0; 0)$
\end{minipage}}

\zadtask{На рисунку зображено квадрат $ABCD$, діагональ якого дорівнює $8\sqrt{2}$~см. Установіть відповідність між відрізком (1--3) та його довжиною (А--Д).
\vspace{0.3cm}
\begin{center}
\begin{tikzpicture}[scale=0.7, thick, line join=round]
    \coordinate (A) at (0,0);
    \coordinate (B) at (0,4);
    \coordinate (C) at (4,4);
    \coordinate (D) at (4,0);
    \draw (A) -- (B) -- (C) -- (D) -- cycle;
    \node[below left] at (A) {$A$};
    \node[above left] at (B) {$B$};
    \node[above right] at (C) {$C$};
    \node[below right] at (D) {$D$};
\end{tikzpicture}
\end{center}
\vspace{0.2cm}
\noindent
\begin{minipage}[t]{0.50\textwidth}
\textit{Відрізок}\par\vspace{0.15cm}
\textbf{1} \quad відстань між серединами сторін $AB$ і $AD$\\[0.25cm]
\textbf{2} \quad радіус кола, уписаного в квадрат $ABCD$\\[0.25cm]
\textbf{3} \quad медіана трикутника $ABC$, проведена з вершини $A$
\end{minipage}
\hfill
\begin{minipage}[t]{0.28\textwidth}
\textit{Довжина відрізка}\par\vspace{0.15cm}
\textbf{А} \quad $2\sqrt{2}$~см\\[0.15cm]
\textbf{Б} \quad $4$~см\\[0.15cm]
\textbf{В} \quad $4\sqrt{2}$~см\\[0.15cm]
\textbf{Г} \quad $8$~см\\[0.15cm]
\textbf{Д} \quad $4\sqrt{5}$~см
\end{minipage}
\hfill
\begin{minipage}[t]{0.18\textwidth}
\vspace{0.4cm}
\begin{flushright}\matchingGrid\end{flushright}
\end{minipage}}
\newpage
\instructionBox{Розв'яжіть завдання 19--22. Одержані числові відповіді запишіть у спеціально відведеному місці.}

\zadtask{Знайдіть площу фігури, обмеженої графіком функції $y=3-\sqrt{x}$ та осями координат.}
\nmtAnswerBox

\zadtask{У будинку $150$ квартир: однокімнатні, двокімнатних на $20\,\%$ менше, ніж однокімнатних, а трикімнатних~--- стільки ж, скільки двокімнатних. Скільки в будинку двокімнатних квартир?}
\nmtAnswerBox

\noindent

\begin{minipage}[c]{0.50\textwidth}
\zadnum На рисунку зображено пряму чотирикутну призму й чотирикутну піраміду $SABCD$, основою якої є ромб $ABCD$. Основою призми є прямокутник, більша сторона якого дорівнює $20$~см, а менша сторона дорівнює діагоналі $AC$ ромба. $BD = 20$~см, $SA = SC = 25$~см, $SB = SD = 26$~см. Визначте об'єм (у см$^3$) призми, якщо площа її бічної поверхні дорівнює $680$~см$^2$.
\end{minipage}%
\hfill
\begin{minipage}[c]{0.48\textwidth}
\centering
\begin{tikzpicture}[scale=0.6, thick, line join=round]
    % ---- Прямокутна призма (зліва) ----
    \def\dx{0.85}\def\dy{0.6}
    \coordinate (F1) at (0,0);
    \coordinate (F2) at (2.2,0);
    \coordinate (F3) at (2.2,3.4);
    \coordinate (F4) at (0,3.4);
    \coordinate (K1) at (\dx,\dy);
    \coordinate (K2) at ({2.2+\dx},\dy);
    \coordinate (K3) at ({2.2+\dx},{3.4+\dy});
    \coordinate (K4) at (\dx,{3.4+\dy});
    \draw (F1)--(F2)--(F3)--(F4)--cycle;     % передня грань
    \draw (F4)--(K4)--(K3)--(F3);            % верхня грань
    \draw (F2)--(K2)--(K3);                   % права грань
    \draw[dashed] (F1)--(K1);
    \draw[dashed] (K1)--(K2);
    \draw[dashed] (K1)--(K4);

    % ---- Піраміда SABCD (справа) ----
    \begin{scope}[shift={(4.6,0.2)}]
        \coordinate (A) at (0,0);
        \coordinate (D) at (3.2,0);
        \coordinate (B) at (1.1,0.9);
        \coordinate (C) at (4.3,0.9);
        \coordinate (S) at (2.15,4.4);
        \coordinate (O) at (2.15,0.45);
        \draw (A) -- (D) -- (C) -- (S) -- cycle;
        \draw (S) -- (D);
        \draw[dashed] (A) -- (B) -- (C);
        \draw[dashed] (S) -- (B);
        \draw[dashed] (A) -- (C);
        \draw[dashed] (B) -- (D);
        \draw[dashed] (S) -- (O);
        \node[below left] at (A) {$A$};
        \node[below right] at (D) {$D$};
        \node[above left] at (B) {$B$};
        \node[above right] at (C) {$C$};
        \node[above] at (S) {$S$};
    \end{scope}
\end{tikzpicture}
\end{minipage}
\par\vspace{0.2cm}
\nmtAnswerBox

\zadtask{Визначте кількість усіх цілих значень $a$, за кожного з яких рівняння
\[
\sqrt{2x - a} \cdot (\log_3(x - 4) - 2) = 0
\]
має два різні корені.}
\nmtAnswerBox

\end{document}